\documentclass[10pt,a4paper]{extarticle}
\usepackage{parskip}

\usepackage[utf8]{inputenc}
\usepackage[T1]{fontenc}
\usepackage[brazil]{babel}
\usepackage{amsmath, amssymb}
\usepackage{geometry}
\usepackage{xcolor}
\usepackage{hyperref}
\usepackage{enumitem}
\usepackage{microtype}
\usepackage{csquotes}
\usepackage{natbib}

\definecolor{important}{RGB}{255, 100, 100}
\definecolor{notice}{RGB}{0, 102, 204}

\newcommand{\impt}[1]{\textcolor{important}{#1}}
\newcommand{\ntc}[1]{\textcolor{notice}{#1}}

\setlist{nosep}
\geometry{margin=2.5cm}

\hypersetup{
  colorlinks=true,
  linkcolor=notice,
  citecolor=notice,
  urlcolor=notice
}

\title{Estudo Dirigido I -- Indutivismo\footnote{Questões baseadas em \citep{chalmers1993}.}}
\author{BCC502 -- Metodologia Científica para Ciência da Computação}
\date{\today}

\begin{document}
\maketitle

\section*{Introdução}
\begin{enumerate}[label=\textbf{\arabic*.}, itemsep=6pt]
  \item Em que consiste o método empírico da física?
  \item O que a filosofia da ciência moderna vê como problema na ideia de ciência baseada em experimento e observação?
\end{enumerate}

\section*{Capítulo 1}
\begin{enumerate}[resume, label=\textbf{\arabic*.}, itemsep=6pt]
  \item Qual é a concepção de senso comum da ciência apresentada por Chalmers?
  \item Qual é a definição de \enquote{indutivismo ingênuo} apresentada por Chalmers?
  \item O que são \enquote{afirmações singulares}?
  \item O que são \enquote{proposições de afirmação}?
  \item O que são \enquote{afirmações universais}?
  \item De acordo com o indutivismo, quais condições devem ser satisfeitas para que afirmações universais possam ser consideradas legítimas?
  \item Dados os conceitos definidos nas questões anteriores, como podemos definir o \enquote{raciocínio indutivo}?
  \item Como é construído o corpo do conhecimento científico de acordo com o indutivista ingênuo?
  \item Para que serve o conhecimento científico?
  \item Explique o processo de dedução.
  \item Apresente de forma esquemática o argumento indutivista da ciência.
  \item Qual é a forma geral de todas as explicações e previsões científicas?
\end{enumerate}

\section*{Capítulo 2}
\begin{enumerate}[resume, label=\textbf{\arabic*.}, itemsep=6pt]
  \item O princípio da indução pode ser justificado? Por quê?
  \item Por que argumentos indutivos não são argumentos logicamente válidos?
  \item Como o indutivista ingênuo pode tentar justificar o princípio da indução? Por que essa justificativa não é válida?
  \item Discuta o \enquote{recuo para a probabilidade} como defesa para o indutivismo. Ele resolve o problema?
  \item Comente sobre as dificuldades relacionadas às condições para generalizações de afirmações singulares para afirmações universais apresentadas no capítulo anterior.
  \item Quais são as três respostas ao problema da indução apresentadas por Chalmers?
\end{enumerate}

\bibliographystyle{apalike}
\bibliography{ref}
\end{document}